%%
%% This is file `ic-max-exemplo.tex',
%% generated with the docstrip utility.
%%
%% The original source files were:
%%
%% ufrj-ic.dtx  (with options: `icmaxexemplo')
%% 
%% This is a sample dissertation or thesis which illustrates the use of the `ufrj'
%% document class with the `ufrj-ic' unit style. It is generated from
%% ufrj-ic.dtx by docstrip.
%% 
%% Copyright (C) 2008-2026 Vicente Helano F. Batista, George O. Ainsworth Jr.,
%% Geraldo Xexéo and Eduardo Mangeli. Released under the GNU General Public
%% License version 3 (see the COPYING.txt file).
%% 
\documentclass[dsc]{ufrj}
\usepackage{ufrj-ic}

\usepackage{makeidx}

\addbibresource{exemplo.bib}

\makelosymbols
\makeloabbreviations
\makelosiglas
\makeglossarylist
\makeindex

\defineacronym{ufrj}{UFRJ}{Universidade Federal do Rio de Janeiro}
\defineacronym{ppgi}{PPGI}{Programa de Pós-Graduação em Informática}

\begin{document}
  \title{Título da Tese}
  \foreigntitle{Title of the Thesis}
  \subtitle{subtítulo, quando houver}
  \foreignsubtitle{subtitle, if any}
  \author{Nome do Autor}{Sobrenome}
  \advisor{Nome do Orientador}{Sobrenome}{D.Sc.}{UFRJ}
  \coadvisor{Nome do Coorientador}{Sobrenome}{D.Sc.}{UFRJ}
  % A banca e o que se declara aqui, na ordem: o orientador, que a preside
  % (3.1.2.1.3e), e o primeiro.
  \examiner{Nome do Orientador Sobrenome}{D.Sc.}{PPGI/UFRJ}
  \examiner{Nome do Coorientador Sobrenome}{D.Sc.}{PPGI/UFRJ}
  \examiner{Nome do Primeiro Examinador Sobrenome}{D.Sc.}{PPGI/UFRJ}
  \examiner{Nome do Segundo Examinador Sobrenome}{D.Sc.}{UFRJ}
  \examiner{Nome do Terceiro Examinador Sobrenome}{Ph.D.}{UFF}
  \department{PPGI}
  \date{03}{2026}
  \approvaldate{30 de março de 2026}
  \keyword{primeira palavra-chave}
  \keyword{segunda palavra-chave}
  \keyword{terceira palavra-chave}
  \foreignkeyword{first keyword}
  \foreignkeyword{second keyword}
  \foreignkeyword{third keyword}

  % Folha adicional da Coleta CAPES (3.1.2.1.2), obrigatoria em teses e
  % dissertacoes. A ficha catalografica vem de fichacatalografica.sibi.ufrj.br;
  % com o arquivo em maos, descomente a linha \catalogcard.
  \concentrationarea{Ciência da Computação}
  % A linha de pesquisa so sai na folha de rosto com a opcao de classe
  % linhanafolhaderosto; sem ela, o comando apenas guarda o valor.
  \researchline{Engenharia de Software}
  \productiontype{bibliographic}
  \linkedproject{yes}
  \projectname{Nome do Projeto de Pesquisa ao qual o trabalho se vincula}
  \fundingagency{Coordenação de Aperfeiçoamento de Pessoal de Nível
                  Superior}{CAPES}
  % \catalogcard{ficha.pdf}

  \maketitle
  \frontmatter

  \dedication{A quem leu a norma antes de escrever.}

  \epigrafe{``A citação que abre o trabalho, com a mesma formatação da
    citação longa.''}{Autor da epígrafe}

  % O resumo e o abstract trazem as instrucoes do Projeto Jacurutu: troque-as
  % pelo seu texto.
  \begin{abstract}
    Utilize, de preferência, um só parágrafo para descrever os objetivos de
    seu trabalho acadêmico, o método com informações sobre coleta, tratamento
    e análise dos dados, as técnicas de abordagem, os resultados e suas
    principais conclusões. Esse resumo deve conter de 150 a 500 palavras. Deve
    ser evitado o uso de frases negativas, símbolos ou contrações que não
    sejam de uso corrente, fórmulas, equações, diagramas etc. que não sejam
    absolutamente necessários; quando for indispensável, defini-los na
    primeira vez que aparecerem. Texto adaptado do Manual do SiBI da UFRJ pelo
    Projeto Jacurutu.
  \end{abstract}

  \begin{foreignabstract}
    Translate your abstract into the language of international dissemination
    chosen by you and your advisor. Translate the keywords as well.
  \end{foreignabstract}

  % As listas, na ordem da 3.1.2: ilustracoes (figuras, quadros, programas,
  % algoritmos), tabelas, abreviaturas, siglas e simbolos. O sumario e sempre o
  % ultimo, e a classe confere a ordem.
  \listoffigures
  \listofquadros
  \listofprograms
  \listofalgorithms
  \listoftables
  \printloabbreviations
  \printlosiglas
  \printlosymbols
  \tableofcontents

  \mainmatter

  \chapter{Introdução}
  \label{cap:introducao}
  O texto do trabalho começa aqui. O \sigla{ppgi} é um Programa da
  \sigla{ufrj}, e na segunda vez a sigla sai sozinha: o \sigla{ppgi}.

  A citação segue o sistema da classe \cite{m-4211}, e a citação com o autor
  no texto sai como em \textcite{m-4212}. A citação de mais de três linhas vai
  no ambiente próprio:

  \begin{longquote}
    A citação longa entra com recuo de 4 cm da margem esquerda, em corpo
    menor, sem aspas e com espaçamento simples, como pede a 2.5 do Manual.
    Depois dela, o texto continua no parágrafo seguinte.
  \end{longquote}

  \section{Uma seção}
  O símbolo \symbl{$\alpha$}{taxa de aprendizado} entra na lista de
  símbolos na ordem em que aparece, e a abreviatura \abbrev{Prof.}{Professor}
  e o \acron{LLM}{Large Language Model} entram nas suas. O
  \glossaryterm{termo técnico}{a definição do termo, que sai no glossário}
  entra no glossário, e esta palavra\index{palavra} entra no índice
  remissivo.

  \subsection{Uma subseção}
  A equação destacada sai centralizada, com o número entre parênteses à
  direita (2.9 do Manual):
  \begin{equation}
    \label{eq:exemplo}
    E = \sum_{i=1}^{n} (y_i - \hat{y}_i)^2
  \end{equation}
  e a menção a ela usa o número: a Equação~(\ref{eq:exemplo}).

  \chapter{Ilustrações, quadros e tabelas}
  \label{cap:ilustracoes}

  A legenda vem em cima e a fonte embaixo, as duas alinhadas à esquerda, na
  margem da ilustração (2.10 e Figura 1 do Manual).

  \begin{figure}[htbp]
    \caption{Uma figura estreita, com a legenda na largura dela}
    \illustrationwidth{6cm}
    \centering
    \rule{6cm}{3cm}
    \source{Elaborado pelo autor (2026).}
  \end{figure}

  \begin{quadro}[htbp]
    \caption{Um quadro: informação textual, com as quatro bordas}
    \centering
    \begin{tabular}{|l|l|}
      \hline
      Elemento & Onde está \\
      \hline
      Capa & 3.1.1.1 \\
      Folha de rosto & 3.1.2.1.1 \\
      \hline
    \end{tabular}
    \source{Elaborado pelo autor (2026).}
  \end{quadro}

  \begin{table}[htbp]
    \caption{Uma tabela: dado numérico, sem as bordas laterais}
    \centering
    \begin{tabular}{lr}
      \hline
      Curso & Alunos \\
      \hline
      Mestrado em Informática & 100 \\
      Doutorado em Informática & 50 \\
      \hline
    \end{tabular}
    \source{Elaborado pelo autor (2026).}
  \end{table}

  \chapter{Programas e algoritmos}
  \label{cap:programa}
  A listagem sai com o estilo da linguagem, e o flutuante dela tem lista
  própria:

  \begin{python}[caption={Um programa de exemplo},label={prog:exemplo}]
def soma(a, b):
    return a + b
  \end{python}
  \source{Elaborado pelo autor (2026).}

  \begin{algorithm}[htbp]
    \caption{Um algoritmo de exemplo}
    \KwData{uma lista $L$ de números}
    \KwResult{a soma $s$ dos elementos de $L$}
    $s \leftarrow 0$\;
    \ForEach{$x \in L$}{
      $s \leftarrow s + x$\;
    }
    \Return{$s$}\;
  \end{algorithm}
  \source{Elaborado pelo autor (2026).}

  \chapter{Conclusão}
  \label{cap:conclusao}
  A conclusão do trabalho, que retoma o que o Capítulo~\ref{cap:introducao}
  propôs.

  \appendix
  \chapter{Um apêndice}
  O apêndice é texto do próprio autor, complementar ao trabalho.

  \annex
  \chapter{Um anexo}
  O anexo é documento de terceiro, trazido para o trabalho.

  \backmatter
  \printbibliography
  \printglossarylist
  \printindex
\end{document}
%% 
%%
%% End of file `ic-max-exemplo.tex'.
